This study reveals a unique evolutionary strategy in the TP53 C-terminal domain (CTD) of pigs, where selective degeneration of regulatory elements (88\% PTM loss) is coupled with complete conservation of DNA-binding function. Through comprehensive comparative analysis of eleven mammalian species, we demonstrate that pigs achieve cancer resistance not through gene duplication (as in elephants) or functional degradation (as in zebrafish), but via a 'less-is-more' mechanism that simplifies CTD regulation while maintaining core tumor suppressor activity. Our CTD-focused analysis of GC content, selection pressure (dN/dS = 8.49 in pig CTD), and structural properties provides novel insights into the evolutionary plasticity of tumor suppressor genes. The findings suggest that strategic domain degeneration can be an effective evolutionary strategy for cancer resistance, with potential implications for understanding TP53 regulation and therapeutic targeting in human cancers.
